
The tunes \textbf{Leather Britches (Breeches)}, \textbf{Britches Full
  of Stitches}, and \textbf{Lord MacDonald} are thought to be
variations of each other. Lord Macdonald's Reel was published in
\textbf{Neil Gow's 3rd Collection of Strathspey Reels} (1792) and may
be the source tune.\footnote{The facsimile can be found online in
  Google Books on page 9:
  \url{https://www.google.com/books/edition/A\_Third\_Collection\_of\_Strathspey\_Reels\_c/1AnfujL8qjIC?gbpv=1}}
Many believe Sir Alexander MacDonald (1744--1795) wrote the tune
himself, but he may have had nothing to do with the composition and
his name became associated with it because he played and popularized
it.\footnote{\url{https://tunearch.org/wiki/Annotation:Lord\_MacDonald\_(4)}}

Leather Britches most likely emerged in America shortly thereafter and
is found in print in \textbf{Music for the Piano Dulcimer by Robert
  J. Rudisill} (New York and Kentucky,
1859).\footnote{\url{https://tunearch.org/wiki/Annotation:Leather\_Britches}
  I couldn’t verify this so let me know if you find this publication.}
Perhaps more interestingly\footnote{To me} is that it appears on page
227-228 in \textbf{A Tour Through Indiana in 1840: The Diary of John
  Parsons of Petersburg, Virginia} (New York, 1920):
\begin{quote}
When we returned to the first cabin we found the young people already
dancing, having induced the old fiddler to take his station in one
corner, where he played in a most lugubrious fashion the old tune of
``Leather Breeches.''\footnote{This can be found online in Google
  Books on page 227-228:
  \url{https://www.google.com/books/edition/A\_Tour\_Through\_Indiana\_in\_1840/gN5CAAAAIAAJ?gbpv=1}} \footnote{I'm
  impressed that such a spritely tune can be played `lugubriously'.}
\end{quote}

The tune's title may refer to leather britches, green snap beans that
are strung together with a needle and thread to dehydrate as a means
of food preservation.\footnote{I remember seeing this as a child
  visiting Alabama, but I think it was at a historical society rather
  than at my Grandparents' house.} Or Britches Full of Stitches could
refer to actual leather pants stitched together as in the lyrics:

\begin{quote}
  Leather breeches, full of stitches, \\
  Mammy sewed the buttons on.\footnote{tunearch, Samuel Preston Bayard. \textbf{Hill Country
    Tunes; Instrumental Folk Music of Southwestern Pennsylvania} (1945)}
\end{\quote}

As a child I remember thinking that Britches Full of Stitches was a
warning that if you didn't stop whatever you were doing, you would
receive corporal punishment resulting in a sore behind.

---Kitty Steetle    
